% LaTeX Canvas: Chapter Outline
\chapter{Collaboration Mechanics}
\label{ch:collaboration}

\section*{Purpose}
Collaboration is not a vague soft skill. It is a set of repeatable mechanics that make work legible, reviewable, and safe to change. In computing and data science, effective collaboration relies on shared artifacts (docs, issues, pull requests), disciplined communication (clear requests and responses), and predictable workflows (review, merge, release). This chapter teaches novice-friendly practices for working with other people without creating confusion, duplication, or fragile ``tribal knowledge.''

\section*{Learning objectives}
By the end of this chapter, you should be able to:
\begin{enumerate}
\item Explain why collaboration depends on shared artifacts (docs, issues, PRs) rather than private messages.
\item Write project documentation that supports onboarding and reproduction.
\item Participate in code review as an author and reviewer using a respectful, actionable style.
\item Use comment threads to clarify intent, document decisions, and close the loop.
\item Propose and follow a lightweight workflow for tasks, reviews, and merges.
\item Maintain collaboration hygiene: small changes, clear ownership, explicit decisions, and visible status.
\end{enumerate}

\section*{Running theme: make work visible, small, and easy to review}
If teammates can understand what changed, why it changed, and how to verify it, collaboration scales.

\section{A beginner mental model: collaboration surfaces}
\subsection{Where collaboration happens}
\begin{itemize}
\item \textbf{Repository}: code, notebooks, documentation.
\item \textbf{Issues}: tasks, bugs, questions, decision records.
\item \textbf{Pull requests}: proposed changes + discussion + review.
\item \textbf{Code review comments}: feedback and decisions tied to specific diffs.
\item \textbf{Project board/milestones}: status tracking.
\item \textbf{Chat/meetings}: coordination, but not the system of record.
\end{itemize}

\subsection{The core collaboration loop}
\begin{center}
Plan in issues $\rightarrow$ implement on a branch $\rightarrow$ open PR $\rightarrow$ review/comment $\rightarrow$ revise $\rightarrow$ merge $\rightarrow$ document $\rightarrow$ repeat
\end{center}

\subsection{Roles and responsibilities}
\begin{itemize}
\item \textbf{Author}: proposes changes, provides context, responds to review.
\item \textbf{Reviewer}: protects quality, asks clarifying questions, mentors.
\item \textbf{Maintainer/lead}: sets policy, arbitrates decisions, ensures follow-through.
\end{itemize}

\section{Documentation as collaboration infrastructure}
\subsection{What documentation is for (student framing)}
\begin{itemize}
\item Onboarding: a new collaborator can start quickly.
\item Reproducibility: a collaborator can rerun results.
\item Decision trace: why key choices were made.
\item Operations: how to build/run/test.
\end{itemize}

\subsection{Minimum viable documentation set}
\begin{itemize}
\item \texttt{README.md}: purpose, setup, how to run, outputs.
\item \texttt{CONTRIBUTING.md}: how to contribute, style rules, review expectations.
\item \texttt{CODE\_OF\_CONDUCT.md} (when appropriate): norms and escalation.
\item \texttt{DECISIONS.md} (or ADRs): consequential decisions and rationale.
\item Data dictionary / codebook: dataset meaning and transformations.
\end{itemize}

\subsection{Writing for your audience}
\begin{itemize}
\item Assume the reader is competent but unfamiliar with your local context.
\item Put the most common tasks first.
\item Provide ``copy/paste to run'' commands and expected outputs.
\end{itemize}

\subsection{Documentation upkeep: avoid staleness}
\begin{itemize}
\item Update docs in the same PR as code changes.
\item Treat docs as part of the definition of done.
\item If documentation is missing, file an issue.
\end{itemize}

\section{Work planning mechanics: issues and task decomposition}
\subsection{Why issues beat chat threads}
\begin{itemize}
\item Searchable, linkable, and persistent.
\item Establishes a shared memory.
\item Supports prioritization and assignment.
\end{itemize}

\subsection{Anatomy of a good issue}
\begin{itemize}
\item Clear title: verb + object.
\item Context: why it matters.
\item Definition of done.
\item Evidence: links, examples, errors.
\item Labels: type (bug/task/question), priority, area (data/code/docs).
\end{itemize}

\subsection{Decomposition patterns (novice-friendly)}
\begin{itemize}
\item Split by artifact: data ingestion, cleaning, analysis, visualization, report.
\item Split by risk: do uncertain tasks early.
\item Split by reviewability: tasks that fit in a small PR.
\end{itemize}

\section{Code review fundamentals (what it is and why it works)}
\subsection{Goals of review}
\begin{itemize}
\item Improve correctness, readability, and maintainability.
\item Catch bugs and edge cases.
\item Transfer knowledge across the team.
\item Align changes with project standards.
\end{itemize}

\subsection{What review is not}
\begin{itemize}
\item Not a personal critique.
\item Not a substitute for testing.
\item Not a place to redesign the whole system in one comment.
\end{itemize}

\subsection{Review checklists (student baseline)}
\begin{itemize}
\item Correctness and clarity: does it do what it says?
\item Tests or verification steps: how do we know?
\item Reproducibility: can someone else run it?
\item Documentation: README/comments updated?
\item Data impacts: schema changes, file paths, assumptions documented?
\item Security/privacy: secrets avoided, sensitive data handled properly?
\end{itemize}

\section{Authoring reviewable changes}
\subsection{Small PR discipline}
\begin{itemize}
\item One purpose per PR.
\item Avoid mixing refactors with new features.
\item Keep diffs readable: remove dead code separately.
\end{itemize}

\subsection{PR descriptions that help reviewers}
\begin{itemize}
\item Summary: what changed.
\item Motivation: why.
\item How to test: steps and expected outputs.
\item Screenshots/figures (if relevant).
\item Links: related issues, background docs.
\end{itemize}

\subsection{Pre-review self-check}
\begin{itemize}
\item Re-read your own diff.
\item Run a smoke test.
\item Confirm docs updated.
\item Ensure no secrets or large accidental files are included.
\end{itemize}

\section{Reviewing and commenting mechanics}
\subsection{Comment taxonomy (so feedback is legible)}
\begin{description}
\item[Blocker] Must fix before merge (correctness, security, reproducibility).
\item[Suggestion] Improvement that is optional.
\item[Question] Clarification needed to evaluate.
\item[Nit] Minor style/detail; non-blocking.
\end{description}

\subsection{Writing effective comments}
\begin{itemize}
\item Be specific: point to a line and explain the issue.
\item Provide rationale: tie to a goal (correctness, clarity, consistency).
\item Offer a concrete suggestion when possible.
\item Prefer questions when intent is unclear (see Chapter~\ref{ch:asking-questions} for how to frame a well-structured technical question).
\end{itemize}

\subsection{Tone and collaboration hygiene}
\begin{itemize}
\item Assume good intent.
\item Critique the code, not the person.
\item Use neutral language and avoid sarcasm.
\item Acknowledge good work and improvements.
\end{itemize}

\subsection{Thread management}
\begin{itemize}
\item One issue per thread.
\item Resolve threads when addressed.
\item Summarize decisions in the PR description or decision log when important.
\end{itemize}

\subsection{Responding to reviews as an author}
\begin{itemize}
\item Reply to each thread: explain what you changed or why you disagree.
\item If you disagree, propose an alternative and ask for alignment.
\item Avoid ``drive-by'' changes that introduce new behavior without explanation.
\end{itemize}

\section{Merging, ownership, and handoffs}
\subsection{Definition of done (team agreement)}
\begin{itemize}
\item Review approvals obtained.
\item Tests and smoke checks pass.
\item Documentation updated.
\item Issue linked and closed with a summary.
\end{itemize}

\subsection{Handoff notes}
\begin{itemize}
\item If someone else will continue the work, leave a short status note.
\item Identify next steps and open questions.
\end{itemize}

\section{Collaboration practices for data science artifacts}
\subsection{Notebooks and noisy diffs}
\begin{itemize}
\item Keep notebooks tidy: avoid huge outputs.
\item Clear unnecessary output before merging.
\item Prefer ``notebook as narrative, code in \texttt{src/}''.
\end{itemize}

\subsection{Data and model artifacts}
\begin{itemize}
\item Avoid committing large raw datasets unless policy says otherwise.
\item Document data retrieval and transformations.
\item Store outputs deterministically and link from reports.
\end{itemize}

\subsection{Reproducibility check as a collaboration contract}
\begin{itemize}
\item Another collaborator should be able to recreate the environment and rerun.
\item Include a minimal ``how to reproduce'' section in the PR.
\end{itemize}

\section{Cadences and rituals (lightweight, student-friendly)}
\subsection{Weekly planning}
\begin{itemize}
\item Review open issues, reprioritize, assign.
\item Identify blockers and dependencies.
\end{itemize}

\subsection{Daily/async updates}
\begin{itemize}
\item Short status: what I did, what I will do, what is blocking me.
\item Post updates in issues/PRs when they affect others.
\end{itemize}

\subsection{Retrospectives (optional)}
\begin{itemize}
\item What went well, what was painful, what to change next time.
\item Convert pain points into issues and documentation improvements.
\end{itemize}

\section{Common failure modes and fixes}
\subsection{Work hidden in private messages}
\begin{itemize}
\item Fix: move decisions into issues/PRs; link evidence.
\end{itemize}

\subsection{Mega-PRs that no one can review}
\begin{itemize}
\item Fix: split into smaller PRs; stage refactors separately.
\end{itemize}

\subsection{Review ping-pong and stalled progress}
\begin{itemize}
\item Fix: clarify blockers vs suggestions; propose a decision; time-box debates.
\end{itemize}

\subsection{Unclear ownership}
\begin{itemize}
\item Fix: assign issues; define a driver for each milestone.
\end{itemize}

\subsection{Documentation drift}
\begin{itemize}
\item Fix: require doc updates in PR checklist; treat missing docs as a bug.
\end{itemize}

\section{Worked examples (outline)}
\subsection{Example 1: Turn a chat question into a good issue}
\begin{itemize}
\item Convert vague request into context + definition of done.
\end{itemize}

\subsection{Example 2: Author a reviewable PR}
\begin{itemize}
\item Create a small branch, open PR with testing steps, request review.
\end{itemize}

\subsection{Example 3: Perform a constructive code review}
\begin{itemize}
\item Use the comment taxonomy; request one blocker and one suggestion.
\end{itemize}

\subsection{Example 4: Close the loop after merge}
\begin{itemize}
\item Update docs, close issue with summary, tag next steps.
\end{itemize}

\section{Templates}
\subsection{Template A: PR checklist}
\begin{verbatim}

## PR Checklist

* [ ] Linked issue(s)
* [ ] Summary and motivation included
* [ ] How to test included
* [ ] Docs updated (README/notes)
* [ ] No secrets or large accidental files
* [ ] Smoke test run
  \end{verbatim}

\subsection{Template B: Review comment format}
\begin{verbatim}
Type: Blocker / Suggestion / Question / Nit
What I see:
Why it matters:
Suggested change:
\end{verbatim}

\subsection{Template C: Decision record (lightweight)}
\begin{verbatim}
Decision:
Date:
Context:
Options considered:
Chosen option:
Rationale:
Consequences:
\end{verbatim}

\subsection{Template D: CONTRIBUTING basics}
\begin{verbatim}

* How to set up environment
* Branch naming and PR policy
* Review expectations
* Style/testing expectations
* Where to ask questions
  \end{verbatim}

\section{Exercises}
\begin{enumerate}
\item Write a README for a small class project that a classmate can run.
\item Turn three vague tasks into well-scoped issues with a definition of done.
\item Open a PR with a strong description and testing steps.
\item Review a peer's PR using the taxonomy; request one change and one clarification.
\item Close an issue by summarizing what changed, what remains, and how to reproduce.
\end{enumerate}

\section{One-page checklist}
\begin{itemize}
\item Work is tracked in issues/PRs, not only chat.
\item Documentation exists and is updated with code changes.
\item PRs are small and include ``how to test''.
\item Reviews are respectful, actionable, and categorized.
\item Decisions are recorded and threads are resolved.
\item After merge, issues are closed with a clear summary and links.
\end{itemize}

% \appendix
\section{Quick reference: collaboration norms}
\begin{itemize}
\item Prefer artifacts over memory.
\item Prefer clarity over speed.
\item Prefer small, reversible changes.
\item Ask questions early; document decisions.
\end{itemize}
